\documentclass[10pt,a4paper]{amsart}
\usepackage[margin=19mm]{geometry}
\usepackage{amsmath,amssymb,mathtools}
\usepackage[scr=rsfs,cal=euler]{mathalfa}
\usepackage{xfrac}
\usepackage[unicode,hidelinks,bookmarks=false]{hyperref}
\newcommand{\ie}{i.e.\ }
\newcommand{\cf}{cf.\ }
\newcommand{\resp}{respectively\ }
\newcommand{\Subsection}{\subsection}
\providecommand{\nobreakdash}{\nobreak\hbox{-}}
\setlength{\emergencystretch}{2em}
\multlinegap0pt
\usepackage{float}
\usepackage{amssymb}
\usepackage[shortlabels]{enumitem}
\newtheorem{theorem}{Theorem}
\title{Triod twist cycles and circle rotations}
\author{Sourav Bhattacharya, Ashish Yadav}
\date{Source: arXiv:2512.20147v1 (CC BY 4.0)}
\begin{document}
\maketitle

\noindent\emph{Source and licence.} Triod twist cycles and circle rotations. Version arXiv:2512.20147v1 (CC BY 4.0), distributed by arXiv under the Creative Commons Attribution 4.0 International licence, http://creativecommons.org/licenses/by/4.0/, as recorded in the arXiv licence metadata for this exact version and shown on the abstract page https://arxiv.org/abs/2512.20147v1. Full licence notice: \texttt{licenses/arxiv-2512.20147-CC-BY-4.0.txt}.\par\smallskip

\noindent\textit{Editorial scope.} Counted unit: the article's theorem movement:black:red (source0.tex lines 861--864). The article's complete original proof of it is delivered with it (source0.tex lines 866--925, one \texttt{proof} environment of 60 lines). No mathematics is written, repaired or re-derived: the statement and the proof are the article's own lines, in the article's own notation, and no retained line is substituted or re-flowed.  No further source-local context is needed: every cross-reference of every retained line resolves inside the two delivered spans.  One declared adaptation: exactly 1 ancillary strings inside the retained spans are omitted (image-inclusion commands and, where the source repeats a label, the repeated label definitions) and no other line differs from the source; the figure environments, with their placement, captions and labels, are kept, and the package records the omitted strings in fidelity receipts bound to the affected bodies.  No retained line contains a citation, so the wrapper carries no bibliography and no bibliography entry.  The retained lines contain the article's own diagram code, which the wrapper typesets with the standard tikz packages; no image file is delivered and the package declares no asset dependency.  Editorial formatting only, in addition: an AMS \texttt{amsart} preamble that restates the article's own declarations and macros used by the retained lines, namely the environments \texttt{theorem} on separate counters, together with the standard packages those lines need.  The retained environments print in this document as Theorem 1 in order of appearance, whereas the article prints the same items with the numbering of its own full document.  The complete native source of the article as distributed in the arXiv e-print for this exact version is bundled unchanged as \texttt{sources/191541/source0.tex}.
\begin{theorem}\label{movement:black:red}
	Let $b \in P$ be any black point. Then there exists a red point $r \in P$ such that 
	$	L(r) - L(b) \leqslant 5\rho - \tfrac{5}{3}$,	where $\rho$ is the rotation number of the triod–twist cycle $P$.
\end{theorem}

\begin{proof}
	Since $\rho >  \tfrac{1}{3}$, the cycle $P$ must contain at least one \emph{red} point.  
	Choose any \emph{red point} $r \in P$, and assume without loss of generality that $r \in b_0$, where $b_0, b_1, b_2$ denote the \emph{branches} of $Y$ in canonical order.  
	We now determine the desired bound by examining the possible positions of the \emph{black point} $b$.
	
	\smallskip 
\smallskip 
\smallskip 

Case 1: $b \in b_0$.  
If $b \geqslant r$, then clearly $L(r) - L(b) \leqslant 0$, so we may assume that $r \geqslant b$. Observe that the sets $\{t \in P : t \geqslant r\}$ and $\{t \in P : f(b) \geqslant t\}$ are not invariant under $f$. Consequently, one of the following sub-cases must occur: either there exists a \emph{red point} $\eta \in P$ such that $f(b) \geqslant \eta$ and $f(\eta) \geqslant r$, or there exist \emph{black points} $\xi_1, \xi_2 \in P$ satisfying $f(b) \geqslant \xi_1$, $f(\xi_1) \geqslant \xi_2$, and $f(\xi_2) \geqslant r$.  Since $P$ is \emph{regular}, it cannot \emph{force} a \emph{primitive cycle} of period $2$, and hence the first sub-case is ruled out. Moreover, for the same reason, in the second sub-case we must necessarily have $f(\xi_1) \geqslant \xi_2 > f(r)$.  By the definition of the \emph{code function}, we obtain,  $L(r) - L(b) \leqslant L(f(\xi_2))  - L(b) + \rho - \frac{1}{3} \leqslant  L(\xi_2) -  L(\xi_1) +2 \rho - \tfrac{2}{3} \leqslant  L(\xi_2) -  L(f(\xi_1)) +3\rho - 1 \leqslant 3 \rho -1$. 

	\smallskip 
	\smallskip 
	\smallskip 

	
	Case 2: $b \in b_1$.  	If $f(b) \geqslant f(r)$, then $L(r) - L(b) = L(f(r)) - \rho + \tfrac{2}{3} - \bigl( L(f(b)) - \rho + \tfrac{1}{3} \bigr)	< \tfrac{1}{3}$.  Hence assume $f(r) \geqslant f(b)$.  	Since the sets $\{t \in P : b \geqslant t\}$ and $\{t \in P : t \geqslant r\}$ are not invariant under $f$, one of the following subcases must occur.
	
	
	\emph{Subcase 2.1:} There exist \emph{black points} $\xi_1, \xi_2, \xi_3, \xi_4 \in P$ with  
	$f(b) \geqslant \xi_1$, $f(\xi_1) \geqslant \xi_2$, $f(\xi_2) \geqslant \xi_3$, $f(\xi_3) \geqslant \xi_4$, and $f(\xi_4) \geqslant r$.  
	Then, $L(r) - L(b) \leqslant L(r) - \bigl( L(f(b)) - \rho + \tfrac{1}{3} \bigr) \leqslant L(r) - L(\xi_1) + \rho - \tfrac{1}{3} = L(r) - L(f(\xi_1)) + 2\rho - \tfrac{2}{3} \leqslant L(r) - L(f(\xi_2)) + 3\rho - 1 \leqslant L(r) - L(f(\xi_3)) + 4\rho - \tfrac{4}{3} \leqslant L(r) - L(f(\xi_4)) + 5\rho - \tfrac{5}{3} \leqslant 5\rho - \tfrac{5}{3}$. 
	

	\emph{Subcase 2.2:} There exists a \emph{red point} $\xi$ with  
	$b \geqslant \xi$ and $f(\xi) \geqslant r$.  Then	$L(r) - L(b)	\leqslant L(r) - L(\xi) 	= L(r) - \bigl( L(f(\xi)) - \rho + \tfrac{2}{3} \bigr) 	\leqslant \rho - \tfrac{2}{3}$. 
	
	\smallskip 
\smallskip 
\smallskip 


	Case 3: $b \in b_2$.  	If $b \geqslant f(r)$, then $L(f(r)) - L(b) \leqslant 0$, which yields, $L(r) + \rho - \frac{2}{3} - L(b) \leqslant 0$ and hence, $L(r) - L(b) \leqslant \frac{2}{3} - \rho $.  So assume $f(r) > b$.  Also, since , $P$ is \emph{regular}, we have $ r > f(b)$. Since neither $\{t \in P : b > t\}$ nor $\{t \in P : t > r\}$ is invariant, one of the following holds.
	
		\emph{Subcase 3.1:} There exist \emph{black points} $\xi_1, \xi_2, \xi_3 \in P$ with  $f(b) \geqslant \xi_1$, 
	$f(\xi_1) \geqslant \xi_2$, $f(\xi_2) \geqslant  \xi_3$, and $f(\xi_3)\geqslant r$ (See Figure \ref{drawing r10}). 	Then $	L(r) - L(b) \leqslant L(r) - \bigl( L(f(b)) - \rho + \tfrac{1}{3} \bigr) \le L(r) - L(\xi_1) + \rho - \tfrac{1}{3} = L(r) - L(f(\xi_1)) + 2\rho - \tfrac{2}{3} \le L(r) - L(f(\xi_2)) + 3\rho - 1 \leqslant L(r) - L(f(\xi_3)) + 4\rho - \tfrac{4}{3} \leqslant 4\rho - \tfrac{4}{3}$. 


\begin{figure}[H]
	\centering
	
	\caption{}
	\label{drawing r10}
\end{figure}


	\emph{Subcase 3.2:} There exist a \emph{black point} $\xi$ and a \emph{red point} $\eta$ with  
$f(b) \geqslant \xi$, $f(\xi) \geqslant  \eta$, and $f(\eta)\geqslant r$.  Then $L(r) - L(b)= L(r) - \bigl( L(f(b)) - \rho + \tfrac{1}{3} \bigr) 	\leqslant L(r) - L(\xi) + \rho - \tfrac{1}{3} \leqslant L(r) - L(\eta) + 2\rho - \tfrac{2}{3} =  L(r) - L(f(\eta)) + 3 \rho - \frac{4}{3} \leqslant 3\rho - \tfrac{4}{3}$. 
	


	
	\smallskip 
	\smallskip 
	\smallskip 
 

	Among all cases, the largest bound arises in Subcase~2.1.  Thus, for every \emph{black point} $b \in P$, there is a \emph{red point} $r \in P$ with	$L(r) - L(b) \leqslant 5\rho - \tfrac{5}{3}$. This completes the proof.
\end{proof}

\end{document}
